\ifx\litemode\undefined\documentclass{article}\fi
\usepackage[margin=3cm]{geometry}
\usepackage{amsmath,amssymb}
\usepackage{mathtools}
\usepackage{algorithm}
\usepackage[noend]{algpseudocode}
\usepackage{amsthm}
\usepackage{boxedminipage}
\usepackage{float}
\usepackage{hyperref}
\usepackage{cleveref}

\algrenewcommand\textproc{}

% Custom commands
\newcommand{\target}{\mathsf{target}}
\newcommand{\kind}{\mathsf{kind}}
\newcommand{\height}{\mathsf{height}}
\newcommand{\validator}{\mathsf{validator}}
\newcommand{\justify}{\mathsf{justify}}
\newcommand{\notarize}{\mathsf{notarize}}
\newcommand{\parent}{\mathsf{parent}}
\newcommand{\true}{\mathsf{true}}
\newcommand{\sentfinalize}{\mathsf{sent\_finalize}}
\newcommand{\sentjustify}{\mathsf{sent\_justify}}
\newcommand{\sentnotarize}{\mathsf{sent\_notarize}}
\newcommand{\getKDeep}{\mathsf{getKDeep}}
\newcommand{\getConfirmed}{\mathsf{getConfirmed}}

% Theorem environments
\theoremstyle{plain}
\newtheorem{theorem}{Theorem}
\newtheorem{lemma}{Lemma}

\theoremstyle{definition}
\newtheorem{definition}[theorem]{Definition}

\theoremstyle{remark}
\newtheorem{remark}[theorem]{Remark}

\title{Simplex with Notarizations}

\IfFileExists{lite_full_setup.tex}{% Shared lite/full setup for protocol-spec .tex files.
%
% Files under full/ that include this snippet compile in two modes:
%   - Full (default): all content rendered.
%   - Lite: when \litemode is defined before \documentclass (typically by a
%     wrapper .tex file at the repo root), lemmas, corollaries, propositions,
%     properties, proofs, remarks, assumptions, and conditions are dropped
%     via the comment package. Algorithms, theorem statements, definitions,
%     and surrounding prose remain.
%
% Place \input of this file in each spec's preamble AFTER all \newtheorem
% declarations (or just before \begin{document}). The exclusions use
% \@ifundefined so it's safe even when a file defines only a subset.
%
% Path resolution: \IfFileExists tries the local copy first (when cwd is
% full/), then falls back to full/lite_full_setup.tex (when cwd is the repo
% root, e.g. while compiling a top-level _lite.tex wrapper).
\newif\iffull
\ifx\litemode\undefined\fulltrue\else\fullfalse\fi
\iffull\else
  \usepackage{comment}
  \makeatletter
  \@ifundefined{lemma}{}{\excludecomment{lemma}}
  \@ifundefined{corollary}{}{\excludecomment{corollary}}
  \@ifundefined{proposition}{}{\excludecomment{proposition}}
  \@ifundefined{property}{}{\excludecomment{property}}
  \@ifundefined{proof}{}{\excludecomment{proof}}
  \@ifundefined{remark}{}{\excludecomment{remark}}
  \@ifundefined{assumption}{}{\excludecomment{assumption}}
  \@ifundefined{condition}{}{\excludecomment{condition}}
  \makeatother
\fi
}{% Shared lite/full setup for protocol-spec .tex files.
%
% Files under full/ that include this snippet compile in two modes:
%   - Full (default): all content rendered.
%   - Lite: when \litemode is defined before \documentclass (typically by a
%     wrapper .tex file at the repo root), lemmas, corollaries, propositions,
%     properties, proofs, remarks, assumptions, and conditions are dropped
%     via the comment package. Algorithms, theorem statements, definitions,
%     and surrounding prose remain.
%
% Place \input of this file in each spec's preamble AFTER all \newtheorem
% declarations (or just before \begin{document}). The exclusions use
% \@ifundefined so it's safe even when a file defines only a subset.
%
% Path resolution: \IfFileExists tries the local copy first (when cwd is
% full/), then falls back to full/lite_full_setup.tex (when cwd is the repo
% root, e.g. while compiling a top-level _lite.tex wrapper).
\newif\iffull
\ifx\litemode\undefined\fulltrue\else\fullfalse\fi
\iffull\else
  \usepackage{comment}
  \makeatletter
  \@ifundefined{lemma}{}{\excludecomment{lemma}}
  \@ifundefined{corollary}{}{\excludecomment{corollary}}
  \@ifundefined{proposition}{}{\excludecomment{proposition}}
  \@ifundefined{property}{}{\excludecomment{property}}
  \@ifundefined{proof}{}{\excludecomment{proof}}
  \@ifundefined{remark}{}{\excludecomment{remark}}
  \@ifundefined{assumption}{}{\excludecomment{assumption}}
  \@ifundefined{condition}{}{\excludecomment{condition}}
  \makeatother
\fi
}

\begin{document}
\maketitle

\section{Model and definitions}

We consider a system of~$n$ validators of which at most~$f$ are adversarial, where $n \geq 3f + 1$.
A \emph{quorum} is any set of at least $2n/3$ validators.
Any two quorums overlap in at least $n/3 + 1 > f$ validators, so their intersection contains at least one honest validator.

\begin{definition}[Block]\label{def:block}
A \emph{block} has a \emph{parent} (another block) and contains a set of votes (Definition~\ref{def:vote}).
The \emph{genesis block}~$B_{\mathrm{genesis}}$ has no parent.
\end{definition}

\begin{definition}[Chain and divergence]\label{def:chain}
A \emph{chain} is a sequence of blocks forming a path from the genesis block, where each block's parent is its predecessor in the sequence.
Since a chain is uniquely determined by its last block, we identify chains with their last block and use the two interchangeably.
Two chains \emph{conflict} if neither is a prefix of the other.
\end{definition}

We write $B \preceq B'$ when $B$ is an ancestor of~$B'$ (or $B = B'$), and $B \sim B'$ when $B$ and $B'$ are \emph{compatible}: $B \preceq B'$ or $B' \preceq B$.
For compatible blocks, we write $B > B'$ to mean $B \succ B'$ (i.e., $B$ is a strict descendant of~$B'$); we say $B$ is \emph{longer} or \emph{greater} than~$B'$.

\begin{definition}[Vote]\label{def:vote}
A \emph{vote} is a signed tuple $(\validator,\; \height,\; \mathsf{slot},\; \target,\; \kind,\; \mathsf{finalize\_height})$ where $\validator$ identifies the signer, $\height$ is a height, $\mathsf{slot}$ is a slot number, $\target$ is a block, $\kind \in \{\notarize, \justify\}$ distinguishes the vote type, and $\mathsf{finalize\_height}$ is either a height or~$\bot$.
Setting $\mathsf{finalize\_height} = h_j$ implicitly commits to finalize the unique block justified at~$h_j$ (Lemma~\ref{lem:uniqueness}), so no separate finalize target is needed.
A vote included on a chain is verified to reference a target that actually exists on that chain.
\end{definition}

\begin{definition}[Fixed-target justification]\label{def:justification}
A block~$T$ is \emph{justified at height~$h$ and slot~$s$} if at least $2n/3$ validators cast $\justify$ votes with $\height = h$, $\mathsf{slot} = s$, and $\target = T$:
\[
  \bigl|\{i \,:\, \mathsf{jtargets}[i] = (T, s)\}\bigr| \;\geq\; 2n/3.
\]
Only votes with matching target \emph{and} slot contribute.
\end{definition}

\begin{definition}[Fixed-target notarization]\label{def:notarization}
A block~$T$ is \emph{notarized at height~$h$ and slot~$s$} if at least $2n/3$ validators cast $\notarize$ votes with $\height = h$, $\mathsf{slot} = s$, and $\target = T$:
\[
  \bigl|\{i \,:\, \mathsf{ntargets}[i] = (T, s)\}\bigr| \;\geq\; 2n/3.
\]
Only $\notarize$ votes contribute.
\end{definition}

\begin{definition}[Finality]\label{def:finality}
The justified block~$J$ is \emph{finalized} when $|P| \geq 2n/3$.
We say $J$ is \emph{finalized at height~$h_j$}.
\end{definition}

\begin{definition}[Slashing conditions]\label{def:slashing}
A validator is \emph{slashable} (or an \emph{equivocator}) if it signed two votes~$a$, $b$ satisfying either condition:

\emph{E1 (double justify target):} both are $\justify$ votes at the same height with different targets:
\[
  a.\kind = \justify
  \;\wedge\;
  b.\kind = \justify
  \;\wedge\;
  a.\height = b.\height
  \;\wedge\;
  a.\target \neq b.\target.
\]

\emph{E2 (finalize--notarize conflict):} $a$ carries a finalize commitment at a height where $b$ is a $\notarize$ vote:
\[
  a.\mathsf{finalize\_height} = b.\height
  \;\wedge\;
  b.\kind = \notarize.
\]

That is: (E1)~a validator cannot cast two $\justify$ votes at the same height with different targets; (E2)~a validator who commits to finalize at height~$h$ (by signing $\mathsf{finalize\_height} = h$) cannot also cast a $\notarize$ vote at height~$h$.
There is no double-target condition for $\notarize$ votes: validators may cast multiple $\notarize$ votes at the same height without slashing risk.
\end{definition}

\section{Algorithm}

\emph{Convention.} Maps accessed at an uninitialized key return $\bot$ by default.

\begin{figure}[H]
\caption{State machine}\label{alg:state-machine}
\begin{center}
\begin{boxedminipage}{\textwidth}
\begin{algorithmic}[1]

\Function{$\textsf{processBlock}$}{$s,\; B$}
  \For{each vote $v$ in $B$} $s \gets \textsf{processVote}(s,\, v,\, B)$ \EndFor
  \State \Return $\textsf{processHeight}(s,\, B)$
\EndFunction
\vspace{3pt}

\Function{$\textsf{processVote}$}{$(h,\, \mathsf{jtargets},\, \mathsf{ntargets},\, J,\, h_j,\, s_j,\, F,\, P,\, N,\, h_n,\, s_n),\; v,\; B$}
  \State $i \gets v.\validator$
  \If{$v.\height = h$ and $v.\target \preceq B$}
    \If{$v.\kind = \justify$ and $\mathsf{jtargets}[i] = \bot$}
      \State $\mathsf{jtargets}[i] \gets (v.\target,\; v.\mathsf{slot})$
    \EndIf
    \If{$v.\kind = \notarize$} \Comment{Only notarize votes}
      \State $\mathsf{ntargets}[i] \gets (v.\target,\; v.\mathsf{slot})$
    \EndIf
  \EndIf
  \If{$v.\mathsf{finalize\_height} = h_j$}
    \State $P \gets P \cup \{i\}$
  \EndIf
  \State \Return $(h,\, \mathsf{jtargets},\, \mathsf{ntargets},\, J,\, h_j,\, s_j,\, F,\, P,\, N,\, h_n,\, s_n)$
\EndFunction
\vspace{3pt}

\Function{$\textsf{processHeight}$}{$(h,\, \mathsf{jtargets},\, \mathsf{ntargets},\, J,\, h_j,\, s_j,\, F,\, P,\, N,\, h_n,\, s_n),\; B$}
  \If{$|P| \geq 2n/3$} \Comment{Finality}
    \State $F \gets J$
  \EndIf
  \If{$\exists\, (T, s)$ s.t.\ $|\{i : \mathsf{jtargets}[i] = (T, s)\}| \geq 2n/3$} \Comment{Justification}
    \If{$T \succeq N$} \Comment{Notarized block advances (justification is a notarization)}
      \State $J \gets T$, \; $h_j \gets h$, \; $s_j \gets s$, \; $P \gets \emptyset$, \; $N \gets T$, \; $h_n \gets h$, \; $s_n \gets s$
    \EndIf
    \State $h \gets h + 1$
    \For{all $i$} $\mathsf{jtargets}[i] \gets \bot$, \; $\mathsf{ntargets}[i] \gets \bot$ \EndFor
    \State \Return $(h,\, \mathsf{jtargets},\, \mathsf{ntargets},\, J,\, h_j,\, s_j,\, F,\, P,\, N,\, h_n,\, s_n)$
  \EndIf
  \If{$\exists\, (T', s)$ s.t.\ $|\{i : \mathsf{ntargets}[i] = (T', s)\}| \geq 2n/3$} \Comment{Notarization}
    \If{$T' \succeq N$}
      \State $N \gets T'$, \; $h_n \gets h$, \; $s_n \gets s$
    \EndIf
    \State $h \gets h + 1$
    \For{all $i$} $\mathsf{jtargets}[i] \gets \bot$, \; $\mathsf{ntargets}[i] \gets \bot$ \EndFor
    \State \Return $(h,\, \mathsf{jtargets},\, \mathsf{ntargets},\, J,\, h_j,\, s_j,\, F,\, P,\, N,\, h_n,\, s_n)$
  \EndIf
  \State \Return $(h,\, \mathsf{jtargets},\, \mathsf{ntargets},\, J,\, h_j,\, s_j,\, F,\, P,\, N,\, h_n,\, s_n)$
\EndFunction

\end{algorithmic}
\end{boxedminipage}
\end{center}
\end{figure}

Height advances via one of two mechanisms, checked in order: (1)~fixed-target justification ($\geq 2n/3$ $\justify$ votes for the same block), or (2)~fixed-target notarization ($\geq 2n/3$ $\notarize$ votes for the same block).
Only justification updates~$J$, $h_j$, and~$P$; notarization advances the height without changing the justified or finalized checkpoints.
The state after block~$B$ is written $\sigma[B]$; the genesis state is $\sigma[B_{\mathrm{genesis}}] = (0,\, \bot^n,\, \bot^n,\, B_{\mathrm{genesis}},\, 0,\, 0,\, B_{\mathrm{genesis}},\, \emptyset,\, B_{\mathrm{genesis}},\, 0,\, 0)$.
Entries in $\mathsf{jtargets}$ and $\mathsf{ntargets}$ are pairs $(\target,\, \mathsf{slot})$; justification and notarization require $2n/3$ votes with matching target \emph{and} slot.

\section{Accountable safety}

\begin{lemma}[Height progression]\label{lem:heightprog}
$\textsf{processHeight}$ increments~$h$ by at most~$1$ (via justification or notarization) or leaves it unchanged.
For any chain reaching height~$h' > 0$, every intermediate height was advanced by some quorum.
\end{lemma}

\begin{lemma}[Justification uniqueness per height]\label{lem:uniqueness}
Unless $\geq n/3$ validators are slashable: at most one block can be justified at any given height across all chains.
If $T_1$ is justified at height~$h$ on chain~$X$ and $T_2$ at height~$h$ on chain~$Y$, then $T_1 = T_2$.
\end{lemma}

\begin{proof}
Justification of~$T_1$ requires $Q_1$ ($\geq 2n/3$) each with $\mathsf{jtargets}[i] = T_1$.
Justification of~$T_2$ requires $Q_2$ ($\geq 2n/3$) each with $\mathsf{jtargets}[i] = T_2$.
By quorum intersection: $|Q_1 \cap Q_2| \geq n/3 > f$, so at least one honest validator is in both.
An honest validator casts at most one $\justify$ vote per height (condition E1), so $T_1 = T_2$.
\end{proof}

\begin{lemma}[Finalization excludes notarization]\label{lem:finexclude-notarize}
Unless $\geq n/3$ validators are slashable: if $C$ is finalized at height~$h$, then no chain can advance past height~$h$ via notarization.
\end{lemma}

\begin{proof}
Finalization of~$C$ at~$h$ requires $Q_F$ ($\geq 2n/3$) each with $\mathsf{finalize\_height} = h$.
Notarization at~$h$ requires $Q_N$ ($\geq 2n/3$) each with a $\notarize$ vote at~$h$.
Intersection $\geq n/3$: each member signed a vote with $\mathsf{finalize\_height} = h$ and a $\notarize$ vote at~$h$, satisfying E2.
\end{proof}

\begin{lemma}[Any chain past a finalized height has $J \succeq C$]\label{lem:mainsafety}
Unless $\geq n/3$ validators are slashable: if $C$ is finalized at height~$h$, then any chain~$Y$ that progresses past height~$h$ has $J \succeq C$ at all heights $> h$.
\end{lemma}

\begin{proof}
Chain~$Y$ advanced through height~$h$ via justification or notarization.

\emph{Notarization:} excluded by \Cref{lem:finexclude-notarize}.

\emph{Justification:} some block~$T$ was justified at~$h$ on chain~$Y$.
By \Cref{lem:uniqueness}, $T = C$.
The justification branch sets $J \gets C$ if $C \succ J$; otherwise $J \succeq C$ already.
Either way $J \succeq C$ after height~$h$, and $J$ is non-decreasing, so $J \succeq C$ at all subsequent heights.
\end{proof}

\begin{lemma}[Finalized blocks form a chain]\label{lem:finchain}
Unless $\geq n/3$ validators are slashable: if $(C, h)$ and $(C', h')$ are finalized checkpoints with $h \leq h'$, then $C' \succeq C$.
\end{lemma}

\begin{proof}
If $h = h'$: $C = C'$ by \Cref{lem:uniqueness}.
If $h' > h$: the chain where $C'$ was finalized progressed past height~$h$.
By \Cref{lem:mainsafety}, $J \succeq C$ on that chain.
$C'$ was the justified block when the finalize quorum formed, so $C' = J \succeq C$.
\end{proof}

\begin{theorem}[Accountable safety]\label{thm:safety}
Unless $\geq n/3$ validators are slashable, no two conflicting blocks can be finalized.
\end{theorem}

\begin{proof}
Suppose $C$ finalized at~$h$ and $C'$ at~$h'$ with $h \leq h'$.
By \Cref{lem:finchain}, $C' \succeq C$, so $C \sim C'$.
\end{proof}

\section{Fork-choice store}

\begin{figure}[H]
\caption{Store and voting}\label{alg:store}
\begin{center}
\begin{boxedminipage}{\textwidth}
\begin{algorithmic}[1]

\State \textbf{Store:} $R_s \gets B_{\mathrm{genesis}}$, \; $F_s \gets B_{\mathrm{genesis}}$, \; $\mathcal{T} \gets \{B_{\mathrm{genesis}}\}$, \; $C \gets \emptyset$
\State $\sigma[B_{\mathrm{genesis}}] \gets (0,\; \bot^n,\; \bot^n,\; B_{\mathrm{genesis}},\; 0,\; 0,\; B_{\mathrm{genesis}},\; \emptyset,\; B_{\mathrm{genesis}},\; 0,\; 0)$
\State \textbf{Per-validator:} $\sentjustify \gets \bot^*$, \; $\sentnotarize \gets \bot^*$, \; $\sentfinalize \gets \bot^*$
\vspace{3pt}

\Function{$\textsf{onBlock}$}{$B$}
  \State \textbf{assert} $B.\parent \in \mathcal{T}$ \Comment{Parent already accepted}
  \State \textbf{assert} $F_s \preceq B$ \Comment{$B$ descends from finalized}
  \State $\mathcal{T} \gets \mathcal{T} \cup \{B\}$
  \State $\sigma[B] \gets \textsf{processBlock}(\sigma[B.\parent],\; B)$
  \State $C \gets C \cup \{(\sigma[B].N,\; \sigma[B].h_n,\; \sigma[B].s_n)\}$
  \State $R_s \gets \textsf{computeGetHeadRoot}()$
  \State $\textsf{updateFinalized}(\sigma[B].F)$
\EndFunction
\vspace{3pt}

\Function{$\textsf{computeGetHeadRoot}$}{}
  \State $(T^*,\, h^*,\, s^*) \gets \arg\max_{(T,\, h,\, s) \in C,\; T \succeq F_s} (h,\; s,\; \mathsf{length}(T))$
  \State \Return $T^*$
\EndFunction
\vspace{3pt}

\Function{$\textsf{updateFinalized}$}{$F_B$}
  \If{$F_B \succ F_s$ and $F_B \preceq R_s$}
    \State $F_s \gets F_B$
  \EndIf
\EndFunction
\vspace{3pt}

\Function{$\textsf{onVote}$}{}
  \State $h^\mathsf{can}\gets \sigma[\textsf{getHead}()].h$
  \State $s^\mathsf{can}\gets \textsf{currentSlot}()$
  \State $(J_j, h_j, \_) \gets \arg\max_{(T,h,s)\in C}(h, s, \mathsf{length}(T))$
  \State $\mathsf{fh} \gets \bot$ \Comment{Finalize height to attach}
  \If{$\lnot\, \sentfinalize[h_j] \land \lnot\, \sentnotarize[h_j]$}
    \State $\mathsf{fh} \gets h_j$
    \State $\sentfinalize[h_j]\gets \true$
  \EndIf
  \If{$\lnot\, \sentjustify[h^\mathsf{can}]$}
    \State \textbf{send} $(i, h^\mathsf{can}, s^\mathsf{can}, \getConfirmed(), \justify, \mathsf{fh})$
    \State $\sentjustify[h^\mathsf{can}]\gets \true$
  \ElsIf{$\mathsf{fh} = \bot$ or $\mathsf{fh} \neq h^\mathsf{can}$} \Comment{E2: avoid notarize at finalize height}
    \State \textbf{send} $(i, h^\mathsf{can}, s^\mathsf{can}, \getKDeep(), \notarize, \bot)$
    \State $\sentnotarize[h^\mathsf{can}]\gets \true$
  \EndIf
\EndFunction

\end{algorithmic}
\end{boxedminipage}
\end{center}
\end{figure}

\end{document}
